
\documentclass{resume} % Use the custom resume.cls style

\usepackage[left=0.75in,top=0.6in,right=0.75in,bottom=0.6in]{geometry} % Document margins
\usepackage{revnum}

\usepackage{hyperref}
\hypersetup{
    colorlinks=true,
    linkcolor=blue,
    filecolor=cyan,      
    urlcolor=magenta,
    pdftitle={Overleaf Example},
    pdfpagemode=FullScreen,
    }
    
\newcommand{\tab}[1]{\hspace{.2667\textwidth}\rlap{#1}}
\newcommand{\itab}[1]{\hspace{0em}\rlap{#1}}
% \name{Ebraheem Khaled Farag} % Your name
% \address{733 Weston Park Drive, Powell Ohio} % Your address
% %\address{123 Pleasant Lane \\ City, State 12345} % Your secondary addess (optional)
% \address{(+1)~614~623~4875 \\ ekfarag@asu.edu} % Your phone number and email
% \address{Orcid: 0000-0002-5794-4286}
\begin{document}
\begin{rSection}{Ebraheem Khaled Farag}
 \centering 
Yale University : Department of Astronomy (YCAA) \\
Arizona State University: School of Earth and Space Exploration (SESE)
\\ Los Alamos National Laboratory: XTD Integrated Design \& Assessment (XTD-IDA)  \hfill 
\\ \textbf{(+1)~614~623~4875, \,  ebraheem.farag@yale.edu, \, Orcid: \href{https://orcid.org/0000-0002-5794-4286}{0000-0002-5794-4286}}
%\\ Astrophysics \hfill { School of Earth and Space Exploration}
\end{rSection}

\begin{rSection}{Research Interests}
Evolution of massive stars: supernova progenitors and pair instability stars
\\Multi-messenger signals: stellar neutrinos and black hole initial mass functions
\\Impact of nuclear reaction uncertainties on nucleosynthesis and stellar evolution
\\Atomic opacities and stellar model micro-physics
\\Asteroseismology of white dwarfs and classical variable stars
\end{rSection}




%----------------------------------------------------------------------------------------
%	Position SECTION
%----------------------------------------------------------------------------------------

\begin{rSection}{Positions}

{\bf YCAA Prize Fellow} \hfill {\em September 2024 - present } 
\\ Yale Center for Astronomy \& Astrophysics Post-doctoral Research Fellowship  \hfill { Yale University}

\end{rSection}

%----------------------------------------------------------------------------------------
%	EDUCATION SECTION
%----------------------------------------------------------------------------------------

\begin{rSection}{Education}

{\bf Arizona State University} \hfill {\em August 2019 - August, 2024 (PhD)} 
\\ PhD Astrophysics \hfill { School of Earth and Space Exploration}
%\\ 5th year PhD Candidate \hfill { Overall GPA: 3.95}
%\\ Astrophysics \hfill { School of Earth and Space Exploration}

{\bf The Ohio State University} \hfill {\em  2014 - 2019 (B.S.)} 
\\ Engineering Physics, Mechanical Specialization\hfill { College of Engineering}
%\\ B.S. \hfill { Overall GPA: 3.76}
%\\ Engineering Physics, Mechanical Specialization\hfill { College of Engineering}


\end{rSection}
%----------------------------------------------------------------------------------------
%	TECHNICAL STRENGTHS SECTION
%----------------------------------------------------------------------------------------

% \begin{rSection}{Technical Strengths}

% \begin{tabular}{ @{} >{\bfseries}l @{\hspace{6ex}} l }
% Computer Languages &  Fortran, MATLAB, Python, C++ \\
% Software & MESA, GYRE, YREC, \LaTeX, Mathematica \\
% \end{tabular}

% \end{rSection}

%----------------------------------------------------------------------------------------
%	WORK EXPERIENCE SECTION
%----------------------------------------------------------------------------------------

% \begin{rSection}{Research Experience}

% \begin{rSubsection}{Los Alamos National Laboratory}{May 2015 - present}{Graduate Student Researcher, Katie Mussack, Suzannah Wood, \& Joyce Guzik}{}
% %\item Developed new solar-evolution-models to explore the solar-abundance-problem. 
% \item Validating proprietary LANL atomic opacity data with modern stellar evolution models and helioseismology.
% %\item Explored the effect of early mass loss, rotation, and magnetic braking on the solar structure.
% \end{rSubsection}

% % \begin{rSubsection}{The Ohio State University}{January 2015 - June 2015}{Student Research Assistant, Department of Physics, Chris Orban}{}
% % \item Used knowledge of high energy-density systems, multiple mean-free-path calculations were computed in MATLAB; to generalize basic constraints necessary for laser shots at the National Ignition Facility (NIF) and OMEGA laser facility.

% % \end{rSubsection}


% \begin{rSubsection}{Arizona State University}{August 2019 - Present}{Graduate Research Assistant, School of Earth and Space Exploration (SESE), Frank Timmes}{}
% \item Using the MESA Stellar evolution code to study gravitational wave detections of merging binary black holes, neutrino astronomy, massive stars, Helium burning nuclear reactions, and Pair-instability supernova.
% \end{rSubsection}









\begin{rSection}{Service}
Modules for Experiment in Stellar Astrophysics (MESA) Software Developer \hfill{\em 2023 - present}
\\Astrophysical Journal Referee - Completed various referee reports for ApJ \hfill {\em  }
\\ SESE Graduate Council Recruitment Chair \hfill {\em 2022-2023} 
\\ Member of Los Alamos National Laboratory Student Association \hfill{\em 2015-2023}
\\ Treasurer -The Society of Physics Students \hfill {2018-2019}
\end{rSection}
\begin{rSection}{Teaching}
TA: Modules for Experiment in Stellar Astrophysics (MESA) Summer School \hfill {\em 2018-2025}
\\ TA: AST 113 Introduction to Astronomy Laboratory\hfill {\em 2020}
\end{rSection}
% \begin{rSubsection}{SESE Graduate Council Recruitment Chair :}{2022-2023}{}
% \end{rSubsection}

% \begin{rSubsection}{Astrophysical Journal Referee  :}{2021}{}
% \item Has completed a referee report for ApJ.\\

\begin{rSection}{Awards}

\begin{rSubsection}{\href{https://aas.org/posts/news/2021/02/chambliss-student-poster-awards-aas-237}{AAS 237 Chambliss Honorable Mention} :}{2021}{}
\item Exploring the Black Hole Mass Gap Sensitivity to $^{12}$C($\alpha$,$\gamma$)$^{16}$O. 
\end{rSubsection}
\end{rSection}


%----------------------------------------------------------------------------------------
\begin{rSection}{Publications} %\itemsep -3pt
Summary: 18 peer-reviewed papers and 5 other scholarly publications/software records; 5 first-author and 3 second-author papers. \\
390 citations; h-index: 10 (\href{https://ui.adsabs.harvard.edu/search/filter_author_facet_hier_fq_author=OR&filter_author_facet_hier_fq_author=author_facet_hier%3A%221%2FFarag%2C%20E%2FFarag%2C%20Ebraheem%22&filter_author_facet_hier_fq_author=author_facet_hier%3A%221%2FFarag%2C%20E%2FFarag%2C%20E%22&fq=%7B!type%3Daqp%20v%3D%24fq_author%7D&fq_author=(author_facet_hier%3A%221%2FFarag%2C%20E%2FFarag%2C%20Ebraheem%22%20OR%20author_facet_hier%3A%221%2FFarag%2C%20E%2FFarag%2C%20E%22)&q=ebraheem%20farag&sort=date%20desc%2C%20bibcode%20desc&p_=0}{ADS}, accessed July 28, 2026). \\
\textbf{Peer Reviewed Publications}\\
First Author Publications:
\begin{revnumerate}
\item \href{https://ui.adsabs.harvard.edu/abs/2026ApJ..1002..172F/abstract}{"Self-Consistent Nonlinear Classical Cepheid Pulsations During Stellar Evolution with MESA"}
\\ \textbf{Farag, E.}, Bellinger, E. P., Mocz, P., Kalici, S., Smolec, R., Kanbur, S., Bettwy, K., \& Lindsay, C. 2026, ApJ, 1002, 172

\item \href{https://ui.adsabs.harvard.edu/abs/2024ApJ...968...56F/abstract}{"An Expanded Set of Los Alamos OPLIB Tables in MESA: Type-1 Rosseland-mean Opacities and Solar Models"}
\\ \textbf{Farag, E.}, Fontes, C. J., Timmes, F. X., Bellinger, E. P., Guzik, J. A., et al. 2024, ApJ, 968, 56

\item \href{https://doi.org/10.3847/1538-4365/ad0787}{"Stellar Neutrino Emission Across The Mass-Metallicity Plane"}
\\ \textbf{Farag, E.}, Timmes, F. X., Chidester, M. T., Anandagoda, S., \& Hartmann, D. H. 2024, ApJS, 270, 5

\item \href{https://doi.org/10.3847/1538-4357/ac8b83}{"Resolving the Peak of the Black Hole Mass Spectrum"}
\\ \textbf{Farag, E.}, Renzo, M., Farmer, R., Chidester, M. T., \& Timmes, F. X. 2022, ApJ, 937, 112

\item \href{https://doi.org/10.3847/1538-4357/ab7f2c}{"On Stellar Evolution in a Neutrino Hertzsprung-Russell Diagram"}
\\ \textbf{Farag, E.}, Timmes, F. X., Taylor, M., Patton, K. M., \& Farmer, R. 2020, ApJ, 893, 133

\end{revnumerate}

Contributing Author Publications:
\begin{revnumerate}
\item \href{https://ui.adsabs.harvard.edu/abs/2026ApJ..1001...41F/abstract}{"Red-Giant Asteroseismology of Low-Mass Population III Stars"}
\\ Ferreira, T., Bellinger, E. P., \textbf{Farag, E.}, \& Lindsay, C. J. 2026, ApJ, 1001, 41

\item \href{https://ui.adsabs.harvard.edu/abs/2026ApJ..1000..147F/abstract}{"Evolution of Low-Mass Population III Stars: Convection, Mass Loss, Nucleosynthesis, and Neutrinos"}
\\ Ferreira, T., Bellinger, E. P., \textbf{Farag, E.}, \& Lindsay, C. J. 2026, ApJ, 1000, 147

\item \href{https://ui.adsabs.harvard.edu/abs/2026ApJ...998L...4S/abstract}{"Evolutionary Tracks and Spectral Properties of Quasi-stars and Their Correlation with Little Red Dots"}
\\ Santarelli, A. D., \textbf{Farag, E.}, Bellinger, E. P., Natarajan, P., Naidu, R. P., Campbell, C. B., \& Caplan, M. E. 2026, ApJL, 998, L4

\item \href{https://ui.adsabs.harvard.edu/abs/2026ApJ...998..150S/abstract}{"MESA-QUEST: Tracing the Formation of Direct-collapse Black Hole Seeds via Quasi-stars"}
\\ Santarelli, A. D., Campbell, C. B., \textbf{Farag, E.}, Bellinger, E. P., Natarajan, P., \& Caplan, M. E. 2026, ApJ, 998, 150

\item \href{https://doi.org/10.1051/0004-6361/202557899}{"Massive Stellar Cannibals: How Stellar Mergers Drive Mass Loss in Extremely Massive Stars"}
\\ Roman-Garza, J., Fragos, T., Charbonnel, C., Ramírez-Galeano, L., Kruckow, M., \& \textbf{Farag, E.} 2026, A\&A, 707, A163

\item \href{https://doi.org/10.3847/1538-4357/ae1a78}{"Solar-like Oscillations in Accreting Pre-main-sequence Stars: Insights and Prospects"}
\\ Jorgensen, J., Zwintz, K., \textbf{Farag, E.}, Vorobyov, E. I., \& Steindl, T. 2025, ApJ, 995, 223

\item \href{https://doi.org/10.3847/1538-4365/ade717}{"Nuclear Neural Networks: Emulating Late Burning Stages in Core-collapse Supernova Progenitors"}
\\ Grichener, A., Renzo, M., Kerzendorf, W. E., Farmer, R., de Mink, S. E., Bellinger, E. P., Chan, C. K., Chen, N., \textbf{Farag, E.}, et al. 2025, ApJS, 279, 49

\item \href{https://doi.org/10.3847/1538-4357/ad4546}{"Photons from neutrinos: The gamma ray echo of a supernova neutrino burst"}
\\ Lunardini, C., Loeffler, J., Mukhopadhyay, M., \textbf{Farag, E.}, \& Timmes, F. X. 2024, ApJ, 969, 149

\item \href{https://doi.org/10.3847/1538-4357/ace620}{"Seismic Signatures of the $^{12}$C($\alpha$, $\gamma$)$^{16}$O Reaction Rate in White Dwarf Models with Overshooting"}
\\ Chidester, M. T., Timmes, F. X., \& \textbf{Farag, E.} 2023, ApJ, 954, 51

\item \href{https://doi.org/10.3847/2041-8213/ac9d39}{"JWST Imaging of Earendel, the Extremely Magnified Star at Redshift z = 6.2"}
\\ Welch, B., Coe, D., Zackrisson, E., de Mink, S. E., Ravindranath, S., Anderson, J., Brammer, G., Bradley, L., Yoon, J., Kelly, P., Diego, J. M., Windhorst, R., Zitrin, A., Dimauro, P., Jimenez-Teja, Y., Abdurro'uf, Nonino, M., Acebron, A., Andrade-Santos, F., Avila, R. J., Bayliss, M. B., Benitez, A., Broadhurst, T., Bhatawdekar, R., Bradac, M., Caminha, G. B., Chen, W., Eldridge, J., \textbf{Farag, E.}, Florian, M., Frye, B., et al. 2022, ApJL, 940, L1

\item \href{https://doi.org/10.3847/1538-4357/ac7ec3}{"On Trapped Modes in Variable White Dwarfs as Probes of the $^{12}$C($\alpha$, $\gamma$)$^{16}$O Reaction Rate"}
\\ Chidester, M. T., \textbf{Farag, E.}, \& Timmes, F. X. 2022, ApJ, 935, 21

\item \href{https://doi.org/10.3847/1538-4357/ac3130}{"Observing Intermediate-mass Black Holes and the Upper Stellar-mass Gap with LIGO and Virgo"}
\\ Mehta, A. K., Buonanno, A., Gair, J., Miller, M. C., \textbf{Farag, E.}, deBoer, R. J., Wiescher, M., \& Timmes, F. X. 2022, ApJ, 924, 39

\item \href{https://doi.org/10.3847/1538-4357/abdec4}{"On the Impact of $^{22}$Ne on the Pulsation Periods of Carbon-Oxygen White Dwarfs with Helium-dominated Atmospheres"}
\\ Chidester, M. T., Timmes, F. X., Schwab, J., Townsend, R. H. D., \textbf{Farag, E.}, Thoul, A., Fields, C. E., Bauer, E. B., \& Montgomery, M. H. 2021, ApJ, 910, 24

\end{revnumerate}

\textbf{Other Scholarly Publications and Software}
\begin{revnumerate}
\item \href{https://arxiv.org/abs/2606.21824}{"A Grid of Fast-rotating, Chemically Homogeneous Supernova and/or Long-GRB Progenitors"}
\\ Renzo, M., Gottlieb, O., Chan, H. S., Goldberg, J. A., Grichener, A., Sen, K., Shah, N., \textbf{Farag, E.}, \& Cantiello, M. 2026, arXiv:2606.21824

\item \href{https://ui.adsabs.harvard.edu/abs/2026ascl.soft04006S/abstract}{"MESA-QUEST: MESA Quasi-star Evolutionary Simulation Toolkit"}
\\ Santarelli, A. D., Campbell, C. B., \textbf{Farag, E.}, et al. 2026, Astrophysics Source Code Library, ascl:2604.006

\item \href{https://arxiv.org/abs/2502.07665}{"Cosmic Massive Star Cannibals and Their Detectability with Gravitational-wave Observations"}
\\ Arca Sedda, M., Scolnic, D., Bortolas, E., Goel, A., Gration, N., Sagar, R., \textbf{Farag, E.}, et al. 2025, arXiv:2502.07665

\item \href{https://doi.org/10.48550/arXiv.2002.04073}{"Investigating Opacity Modifications and Reaction Rate Uncertainties to Resolve the Cepheid Mass Discrepancy"}
\\ Guzik, J. A., \textbf{Farag, E.}, Ostrowski, J., et al. 2021, in Astronomical Society of the Pacific Conference Series, Vol. 529, RR Lyrae/Cepheid 2019: Frontiers of Classical Pulsators, ed. K. Kinemuchi, C. Lovekin, H. Neilson, \& K. Vivas, 79

\item \href{https://doi.org/10.1017/S1743921316006062}{"Sound speed and oscillation frequencies for solar models evolved with Los Alamos ATOMIC opacities"}
\\ Guzik, J. A., Fontes, C. J., Walczak, P., et al. 2016, IAU Focus Meeting, 29B, 532

\end{revnumerate}
\end{rSection}

\begin{rSection}{Presentations \& Invited Talks}
\begin{revnumerate}

\item \href{https://www.issibern.ch/}{From Stellar Evolution to Pulsation: Post-Main Sequence Massive Star Pulsations} \hfill {\em Invited Seminar}
\\ International Space Science Institute, Bern, Switzerland, 2026

\item \href{https://wwwmpa.mpa-garching.mpg.de/SESTAS/doku.php?id=start}{Self-consistent Non-linear Radial Pulsations during Stellar Evolution} \hfill {\em Invited Seminar}
\\ Max Planck Institute for Astrophysics, 2025

\item \href{https://astronomy.yale.edu/event/yale-astronomy-colloquium-ebraheem-farag}{Multi-Messenger Signals from Stellar Evolution} \hfill {\em Invited Colloquium}
\\ Yale Astronomy Colloquium, 2025

\item \href{https://ui.adsabs.harvard.edu/abs/2023HEAD...2011520F/abstract}{Neutrino Emission from Stars} \hfill {\em Contributed Poster}
\\ \textbf{Farag, E.}, Chidester, M., \& Timmes, F. 2023, in AAS/High Energy Astrophysics Division, Vol. 55, AAS/High Energy Astrophysics Division, 115.20

\item \href{https://ui.adsabs.harvard.edu/abs/2023HEAD...2011662F/abstract}{Resolving The Peak Of The Black Hole Mass Spectrum} \hfill {\em Contributed Poster}
\\ \textbf{Farag, E.}, Renzo, M., Farmer, R., Chidester, M., \& Timmes, F. 2023, in AAS/High Energy Astrophysics Division, Vol. 55, AAS/High Energy Astrophysics Division, 116.62

\item \href{https://ui.adsabs.harvard.edu/abs/2023AAS...24125904F/abstract}{Resolving The Peak Of The Black Hole Mass Spectrum} \hfill {\em Contributed Talk}
\\ \textbf{Farag, E.}, Renzo, M., Farmer, R., Chidester, M., \& Timmes, F. 2023, in American Astronomical Society Meeting Abstracts, Vol. 55, American Astronomical Society Meeting Abstracts, 259.04

\item \href{https://mesahub.github.io/summer-school-2022/}{Recurrent Novae and Steady Burning} \hfill {\em Invited Lecturer}
\\ (MESA) Summer School, Santa Barbara, CA, 2022

\item \href{https://doi.org/10.2172/1876778}{Resolving The Peak of The Black Hole Mass Spectrum} \hfill {\em Invited Seminar}
\\ \textbf{Farag, E. K.}, Renzo, M., Farmer, R., Chidester, M., \& Timmes, F. 2022
\\ Los Alamos Distinguished Seminar Series, Los Alamos, NM

\item \href{https://ui.adsabs.harvard.edu/abs/2021AAS...23752403F/abstract}{Exploring The Black hole Mass Gap With Revised Nuclear Reaction Rates} \hfill {\em Contributed Poster}
\\ \textbf{Farag, E.}, \& Timmes, F. 2021, in American Astronomical Society Meeting Abstracts, Vol. 53, American Astronomical Society Meeting Abstracts, 524.03

\item \href{https://doi.org/10.2172/1545733}{Is There A Solar Model Solution to The Faint Young Sun Paradox?} \hfill {\em Invited Seminar}
\\ \textbf{Farag, E. K.} 2019
\\ Los Alamos Distinguished Seminar Series, Los Alamos, NM

\item \href{https://ui.adsabs.harvard.edu/abs/2020JAVSO..48..102G/abstract}{Modeling Cepheid Variable Stars Using the Open-Source MESA Code} \hfill {\em Contributed Talk}
\\ Guzik, J. A., \textbf{Farag, E.}, Ostrowski, J., et al. 2020, The Journal of the American Association of Variable Star Observers, vol. 48, no. 1, p. 102

\item \href{https://ui.adsabs.harvard.edu/abs/2020anms.conf....5G/abstract}{Modeling Pulsations of Cepheid Variables using the Open-Source MESA Code} \hfill {\em Contributed Talk}
\\ Guzik, J. A., \textbf{Farag, E.}, Ostrowski, J., et al. 2020, in The 35th Annual New Mexico Symposium, ed. A. D. Kapinska, 5

\item \href{https://ui.adsabs.harvard.edu/abs/2015APS..DPPUP2023V/abstract}{A diffusive radiation hydrodynamics code, xRage, is implemented to compare radiation flow with experimental data from the Omega laser facility} \hfill {\em Contributed Poster}
\\ Vandervort, R., Elgin, L., \textbf{Farag, E.}, et al. 2015, in APS Meeting Abstracts, Vol. 2015, APS Division of Plasma Physics Meeting Abstracts, UP12.023

\end{revnumerate}
\end{rSection}




\begin{rSection}{Accepted Proposals}
\begin{enumerate}
    \item \href{https://ui.adsabs.harvard.edu/abs/2021hst..prop16668C/abstract}{A Strongly Magnified Individual Star and Parsec-Scale Clusters Observed in the First Billion Years at z = 6, JWST Proposal. Cycle 1, ID. \#2282}
    \\ Coe, D., Acebron, A., Avila, R., et al. 2021,  

    \item OzGrav: The Australian Research Council Centre of Excellence for Gravitational Wave Discovery
    \\ For the MESA Down Under Winter School, hosted by the University of Sydney, 2024
\end{enumerate}    
\end{rSection}
\end{document}
